% Neutral guideline modules. No clinical recommendations or source-paper content.
% Shared dependencies and neutral material. Compile with XeLaTeX.
% No publisher logos, paper text, data, figures, or download metadata are included.
\usepackage{amsmath,amssymb,graphicx,xcolor,tikz,fancyhdr,needspace,xurl,hyperref,balance,etoolbox}
% The legacy balance package omits the first split box's final baseline depth.
% Reserve that measured depth rather than hiding a CJK overfull-box warning.
\makeatletter
\patchcmd{\@BAlancecol}{\setbox\@leftcolumn=\vbox to \@tempdima}{\advance\@tempdima by \dp1\relax\setbox\@leftcolumn=\vbox to \@tempdima}{}{\PackageWarning{cns-template}{Balance depth patch unavailable; inspect final columns}}
\makeatother
\usetikzlibrary{arrows.meta,positioning}
\usepackage[protrusion=true,expansion=false]{microtype}
\IfFontExistsTF{Noto Serif CJK SC}{%
  \setCJKmainfont{Noto Serif CJK SC}%
}{%
  \setCJKmainfont[BoldFont=FandolSong-Bold.otf,ItalicFont=FandolKai-Regular.otf]{FandolSong-Regular.otf}%
}
\IfFontExistsTF{Noto Sans CJK SC}{%
  \setCJKsansfont{Noto Sans CJK SC}%
}{%
  \setCJKsansfont[BoldFont=FandolHei-Bold.otf]{FandolHei-Regular.otf}%
}
\setmainfont{TeX Gyre Termes}\setsansfont{TeX Gyre Heros}
\setmonofont{TeX Gyre Cursor}
\definecolor{JournalBlue}{HTML}{0097C9}
\definecolor{JournalRed}{HTML}{B83443}
\definecolor{ScienceRed}{HTML}{CD1021}
\definecolor{CiteBlue}{HTML}{0000EE}
\hypersetup{colorlinks=true,linkcolor=JournalBlue,urlcolor=JournalBlue,
  pdftitle={CNS论文汉化：中性版式示例},pdfauthor={示例模板；非期刊官方模板},
  pdfsubject={Neutral source-inspired Chinese layout demonstration; no real research claims}}
\urlstyle{same}\raggedbottom
\setlength{\parindent}{1em}\setlength{\parskip}{0pt}
\setlength{\emergencystretch}{2em}\tolerance=2200
\XeTeXlinebreaklocale "zh"\XeTeXlinebreakskip=0pt plus .1pt
\pagestyle{fancy}\fancyhf{}
\renewcommand{\headrulewidth}{0pt}\renewcommand{\footrulewidth}{0pt}
\newcommand{\PaperTitle}{中文题名占位：从组织结构到证据解释}
\newcommand{\SampleAuthors}{Author A\textsuperscript{1}, Author B\textsuperscript{2,*}}
\newcommand{\SampleAffiliations}{\textsuperscript{1}作者单位一（示例占位）。\quad
  \textsuperscript{2}作者单位二（示例占位）。\par
  *通信作者信息待填。DOI、收稿和出版日期均应从目标原文核实。}
\newcommand{\SampleAbstract}{本段是中性摘要占位文字，不陈述任何真实研究发现。它演示中文题名、作者信息、摘要和双栏正文之间的层级关系。实际使用时，应完整替换为目标论文的忠实中文译文，并保留研究对象、方法、结果和结论的限定条件。图注、参考文献、作者声明和来源信息均属于独立的覆盖项。模板中的图形仅说明排版结构，不表示生物学机制、实验结果或统计关系。}
\newcommand{\DemoCaption}{\textbf{图1. 中性证据链示意（占位图，无实验数据）。}
  \textbf{A}，输入节点表示需要核对的原文材料。
  \textbf{B}，处理节点表示逐段转换及复核。
  \textbf{C}，输出节点表示正文、图注和书目等独立交付部分。
  箭头仅表示本示例的阅读顺序，不代表因果关系、效应方向或数量变化。
  该图由模板自行绘制，所有节点均为文字占位；无样本量、误差线、统计检验或比例尺。
  实际论文图注应逐项保留原稿报告的样本量、单位、检验方法、限定及图形署名。}
\newcommand{\DemoVector}{%
 \resizebox{\linewidth}{!}{%
 \begin{tikzpicture}[font=\sffamily\small,node distance=18mm,
   box/.style={draw=black!70,rounded corners=2pt,fill=black!3,minimum width=26mm,minimum height=16mm,align=center},
   flow/.style={-{Stealth[length=2.5mm]},line width=.7pt}]
  \node[box] (a) {\textbf{A}\quad 输入\\原文材料};
  \node[box,right=of a] (b) {\textbf{B}\quad 处理\\转换与复核};
  \node[box,right=of b] (c) {\textbf{C}\quad 输出\\中文译稿};
  \draw[flow] (a)--(b);\draw[flow] (b)--(c);
  \node[below=11mm of a,align=center] {标题与摘要\\正文与图表};
  \node[below=11mm of b,align=center] {保留限定\\检查单位};
  \node[below=11mm of c,align=center] {可追溯引文\\明确原文信息};
  \node[below=29mm of b,font=\sffamily\footnotesize,text=black!65] {示意占位；不包含真实实验数据};
 \end{tikzpicture}}%
}
\newcommand{\SampleBibliography}{%
 \section*{参考文献（示例占位）}
 {\fontsize{7.6bp}{10bp}\selectfont\setlength{\parindent}{0pt}\interlinepenalty=10000
 \hypertarget{ref:1}{}1. Author A. (20XX). Example reference title. \textit{Journal placeholder}. 卷、页码与DOI待填。此条仅供演示引文跳转。\par\smallskip
 \hypertarget{ref:2}{}2. Author B. (20YY). Another example reference title. \textit{Journal placeholder}. 出版信息待填；不要把占位条目带入正式译稿。\par}
}
\newcommand{\SampleSourceInfo}{%
 \section*{原文信息}
 以下字段全部为示例占位：作者；英文原标题；期刊；年份；卷期；页码或文章号；DOI或来源链接。实际使用时必须逐项替换并核对。\par
 本文件是\JournalName 风格的非官方中文译稿样板，依据一次已核验文章版式进行结构化改编，不代表该期刊所有年代或文体的统一标准，也不保证像素级复刻。无真实论文内容或期刊标识图像。\par
}
% Deliberately neutral demo prose, not translated from a published paper.
\newcommand{\SampleOpening}{%
 \section*{引言：示例正文与引文}
 这是一段用于检查双栏阅读节奏的中性示例文字。段落中的术语、缩写和标点需要在实际项目中与源文保持一致。中文译稿应呈现原论文的完整论述，而不是把多段证据压缩成一段摘要。这里的引文只链接到本文件末尾明确标为占位的参考文献\SampleCiteOne。\par
 若原句仅表示相关、可能或提示，译文也应保留相同的证据强度。数字、单位、正负号、上下标和范围端点应逐项核对。原文存在内部矛盾时，应分别保留相应位置的陈述，再用简短译注说明；不能根据常识把不同的数字静默统一。\par
 \subsection*{内容层级与跨页衔接}
 此小标题演示二级层级。实际论文可能采用不同的小标题密度，模板不会替代对原文结构的识别。双栏文章需按左栏、右栏和跨页续句的顺序读取。若段落在页末中断，应在完成整句核对后接入下一页，避免把图中文字误当正文，或重复翻译同一段。\par
 图表具有独立的信息范围。正文引用到某幅图，并不意味着其图注可以省略。图注中的每个分面、比较对象、时间点、样本量及统计说明都应对应源文。英文图内标签可以保留，并在图旁或文末提供中文释义；实际采用哪种安排，需要依据图形尺寸和译文长度判断\SampleCiteTwo。\par
 \subsection*{来源与可维护数据}
 本段继续演示正文密度。可编辑文本和图注数据应与最终版面分开维护，以便修订时准确定位变化。若只调整版式，应能证明正文、图注和书目内容没有被意外改写。更换字体或纸张后，须重新检查断行、图注与图片的相对位置以及页脚可读性。\par
 文内引用与书目之间应有明确对应关系。这里的两条演示书目没有真实DOI，也没有虚构的实验信息。使用模板时可增删条目，但必须同时更新编号、链接目标和引用位置。页码只表示中文译稿自己的分页，不用于冒充原出版页码。\par
 \subsection*{结果叙述的演示占位}
 本节的文字仅用于呈现长段落在双栏页面中的流动，不报告真实结果。实际翻译应保留观察对象、实验条件和结局指标之间的关系。例如，原文分别描述体外实验与动物实验时，不能把两类结果合并为同一层面的结论；研究使用的对照、时间顺序和分析范围也必须保留。图表中的比较范围不应扩大为对整个人群或全部疾病阶段的判断。\par
 论述中的转折与限制同样属于正文。若研究先提出一种解释，再讨论另一种可能，译文应沿用这一推进顺序。作者明确写出的不确定性应能在中文中被识别；如果尚未确定某个现象的方向或机制，就不能改成已经建立的因果关系。对于术语近似但研究层级不同的概念，可在首次出现时保留缩写，后文保持统一。\par
 \subsection*{方法信息与可复核性}
 方法部分往往包含密集数字和专有名词。翻译时须区分浓度、总量、频率和处理时长，不能因为中文句式调整而改变参数的归属。试剂名称、软件版本和数据访问标识通常保留原样；上下限、排除规则和显著性标准则需要逐字符核对。这里没有提供任何真实实验参数，避免把模板占位误认为可以直接执行的方法。\par
 表格与补充材料的引用也需要保留原层级。正文对补充图的指向并不意味着模板中已经附有该补充图；若源文件只包含主文，应如实标注所交付的范围。图表编号、附录顺序和参考文献数量都应由目标文章决定，本示例中的结构不构成固定配额，也不要求所有论文采用同样的章节。\par
}
\newcommand{\SampleClosing}{%
 \section*{讨论：版式适配与范围}
 上方矢量图是模板原创的排版占位图，不来自任何论文。图形与完整演示图注被放在同一页，便于读者直接核对分面说明。真实项目应使用经过授权或用户提供的原图素材，并保留图注署名与许可信息。尺寸不足时，可适度缩放或调整栏宽，不以删除图注解决空间问题。\par
 模板仅记录一次已核验文章的可观察设计特征。即使文章出自同一期刊，不同年代、研究论文与综述、校正清样与正式版本也可能采用不同的页眉和引文形式。应以用户提供的实际PDF为准，不能仅凭期刊名称推断配色、字号或引用方式\SampleCiteOne。\par
 \subsection*{演示检查项}
 本样板用于展示可编译性、中文字形、标题层级、跨栏图注、页眉页脚和内部链接。图中没有实际数据，因此不应被当作科学插图引用。替换内容后，还需检查所有页的渲染结果，确认无截字、碰撞、异常空白和遗漏段落。\par
 \SampleBibliography
 \SampleSourceInfo
}


\usepackage{multicol,tabularx,array,colortbl,booktabs}
\setlength{\columnsep}{15.3bp}
\setlength{\parindent}{1em}\setlength{\parskip}{2bp}
\renewcommand{\normalsize}{\fontsize{9.2bp}{12.2bp}\selectfont}
\definecolor{GuideGrey}{HTML}{EDEEEE}
\definecolor{GuideGreen}{HTML}{6AA78B}
\definecolor{GuideEvidence}{HTML}{678FB4}
\definecolor{GuideYellow}{rgb}{1,.853,.4999}
\definecolor{GuideOrange}{rgb}{.9655,.6886,.3759}
\definecolor{GuideHarm}{rgb}{.9539,.4904,.4976}
\definecolor{GuideHarmDeep}{rgb}{.6918,.2195,.1115}
\ctexset{section={format=\sffamily\bfseries\color{GuideRed}\fontsize{12.4bp}{15bp}\selectfont,beforeskip=10bp,afterskip=4bp},
 subsection={format=\sffamily\bfseries\color{GuideRed}\fontsize{10.4bp}{13bp}\selectfont,beforeskip=8bp,afterskip=3bp}}
\fancyhead[L]{\sffamily\fontsize{8bp}{10bp}\selectfont 示例编写组}
\fancyhead[R]{\sffamily\fontsize{8bp}{10bp}\selectfont 指南文体中性样板}
\renewcommand{\headrulewidth}{.4bp}
\fancyfoot[L]{\sffamily\fontsize{6.7bp}{8.5bp}\selectfont\JournalName｜出版字段按所给版本核实}
\fancyfoot[R]{\sffamily\fontsize{6.7bp}{8.5bp}\selectfont 非官方样板｜中文译页\thepage}
\hypersetup{colorlinks=true,linkcolor=black,urlcolor=GuideRed,
 pdftitle={指南文体中性中文排版样板},pdfauthor={非官方中性样板}}
\newcommand{\GuideCite}[1]{\textsuperscript{\hyperlink{ref:#1}{#1}}}
\newcommand{\GuideBody}{%
本段是原创中性排版文字，不包含真实诊疗建议。指南翻译需要分别核对解释性正文、定义、推荐表、推荐类别、证据等级、算法、普通表格及全部脚注。相同数字在不同章节可能属于不同书目分组，引用目标应随原稿的编号体系保存。\GuideCite{1}

复用已经核对的译文时，应把内容与样式分开维护，保存段落、推荐标记和书目的内容指纹。纸张、字体、色彩、栏结构或分页的变化不应改变否定、条件、范围、推荐强度或证据级别。不能把另一学会的分类体系替换进原表。

图片可按实际大小在通栏或单栏排放，并穿插于正文。图注过长时可在面板或句子边界续页，注明相邻续注；不要把短段单独留在一页，也不要为了填空删掉科学说明。图像保留英文时，还应提供完整的中文标签释义。}
\newcommand{\GuideTableIntro}{%
\noindent{\bfseries\color{GuideRed}表1．推荐类别与证据栏的中性排版示例}\par
{\fontsize{7.6bp}{9.2bp}\selectfont 以下行只描述翻译核对工作，不是真实临床推荐。类别和证据符号仅示范排版，不赋予这些示例文字医学含义。}\par\vspace{4bp}}
\newcommand{\AhaGuideTable}{%
\GuideTableIntro
{\fontsize{8.2bp}{10.5bp}\selectfont\setlength{\tabcolsep}{5bp}\renewcommand{\arraystretch}{1.25}
\noindent\begin{tabularx}{\linewidth}{@{}>{\centering\arraybackslash}p{50bp}>{\centering\arraybackslash}p{42bp}X@{}}
\rowcolor{black!20}\textbf{COR}&\textbf{LOE}&\textbf{推荐内容（中性占位）}\\
\cellcolor{GuideGreen}1&\cellcolor{GuideEvidence}B-R&逐项保留来源中的条件、对象及证据标记，不把此行用于医疗决策。\\
\cellcolor{GuideYellow}2a&\cellcolor{GuideEvidence!60}B-NR&在有数据支持的情况下逐条核对；本行仅展示多行文字与色块对齐。\\
\cellcolor{GuideOrange}2b&\cellcolor{GuideEvidence!35}C-LD&来源不完整时明确说明，不推造缺失等级或改变“可考虑”的强度。\\
\cellcolor{GuideHarm}3（无获益）&\cellcolor{GuideEvidence!20}C-EO&保留来源中“无获益”的原始类别，不把它改写成“有害”。\\
\cellcolor{GuideHarmDeep}\color{white}3（有害）&\cellcolor{GuideEvidence!20}C-EO&有害类别独立显示；本表只示范色彩、换行及类别关系。\\
\end{tabularx}\par}
{\fontsize{7.6bp}{9.2bp}\selectfont COR：推荐类别；LOE：证据等级。ACC/AHA的A、B-R、B-NR、C-LD、C-EO按来源保留；表中五行是原创示例，不是固定推荐数量。}\par}
\newcommand{\EscGuideTable}{%
\GuideTableIntro
{\fontsize{8.2bp}{10.5bp}\selectfont\setlength{\tabcolsep}{5bp}\renewcommand{\arraystretch}{1.25}
\noindent\begin{tabularx}{\linewidth}{@{}X>{\centering\arraybackslash}p{38bp}>{\centering\arraybackslash}p{38bp}@{}}
\rowcolor{black!20}\textbf{推荐内容（中性占位）}&\textbf{Class}&\textbf{LOE}\\
\rowcolor{GuideGrey}完整保留每条推荐的内容和脚注，不把此示例当作诊疗建议。&\cellcolor{GuideGreen}I&\cellcolor{GuideEvidence}A\\
\rowcolor{GuideGrey}根据来源检查限定词和比较方向，等级字段独立核验。&\cellcolor{GuideYellow}IIa&\cellcolor{GuideEvidence!65}B\\
\rowcolor{GuideGrey}原文缺少字段时明确标注，不能用其他体系补齐。&\cellcolor{GuideOrange}IIb&\cellcolor{GuideEvidence!35}C\\
\rowcolor{GuideGrey}否定性推荐保留原意，色彩只帮助识别对应类别。&\cellcolor{GuideHarm}III&\cellcolor{GuideEvidence!35}C\\
\end{tabularx}\par}
{\fontsize{7.6bp}{9.2bp}\selectfont 本表只示范ESC的I、IIa、IIb、III与A、B、C列。真实指南的续表应重复标题、列头及分组关系，并完整保留脚注。}\par}
\newcommand{\GuideFigure}{%
\par\vspace{6bp}\begin{center}\begin{minipage}{.86\linewidth}\DemoVector\end{minipage}\end{center}
{\fontsize{7.6bp}{9.2bp}\selectfont\textbf{\color{GuideRed}图1．内容复用与排版核对（原创占位）。}
\textbf{A}，输入已核对译文、推荐字段及英文书目。
\textbf{B}，更换原刊纸张、题名、栏结构和图注样式，保留内容顺序。
\textbf{C}，编译并逐页检查，验证内容指纹与引文跳转。
箭头只表示本示例的处理顺序，不代表临床路径、诊断或治疗关系。本图无真实数据、患者信息或原刊图形；真实指南须按各图面板完整核对。}\par\vspace{8bp}}
\newcommand{\GuideEnd}{%
\section*{参考文献（中性占位）}
{\fontsize{7.6bp}{9.2bp}\selectfont\setlength{\parindent}{0pt}
\hypertarget{ref:1}{}1. Example author. Example source. All publication fields are placeholders. 非真实文献。\par}
\section*{声明与范围}
推荐示例仅供排版演示，不构成医疗建议。正文、图、表及作者信息均为原创中性占位；正式稿应以用户提供版本为事实来源。未取得的补充材料不能标为已翻译。
\section*{原文信息}
作者、英文原题、期刊、年份、卷期、页码或文章号、DOI及明确出版日期均待按目标文件替换。当前文件为\JournalName 风格的非官方样板，无真实论文全文、原图或下载信息。

仅供学术交流；出版物版权归原作者及出版商所有。}
